\section{Immutable Parent Pointers}
\label{sec:immutable-parent-pointers}

The immutable parent pointer is the foundational structural primitive of the Recursive Spawning Protocol. It is the single mechanism that transforms the delegation chain from a forensic reconstruction attempted after the fact into a runtime-verifiable artifact — one that is immutable, auditable, and architecturally enforced. It renders all modes of autonomous drift structurally impossible regardless of operational urgency, architectural complexity, or adversarial sophistication of any machine actor in the chain.

This section specifies the parent pointer formally (\S\ref{sec:parent-pointer-formal}), defines the chain traversal operator and establishes the base case for the root human anchor (\S\ref{sec:chain-definition}), specifies immutability enforcement through the Fact Layer write protocol (\S\ref{sec:fact-layer-enforcement}), describes the Purpose Layer's spawn authorization role (\S\ref{sec:purpose-layer-authorization}), and provides the complete chain traversal algorithm with correctness arguments (\S\ref{sec:chain-traversal}).

% -------------------------------------------------------------------------- %
\subsection{Formal Definition: ParentPointer}
\label{sec:parent-pointer-formal}

\begin{dfn}[Parent Pointer]
\label{dfn:parent-pointer}
Let $\mathcal{N}$ be the set of all agent nodes in a \bxthree{} system implementing the Recursive Spawning Protocol, and let $P$ be the set of all Purpose Layer actors. The \emph{parent pointer} $\alpha(n)$ is defined for every node $n \in \mathcal{N}$ as the total function
\begin{equation}
\alpha : \mathcal{N} \rightarrow P \cup \{\bot\}
\label{eq:alpha-function}
\end{equation}
where $\alpha(n) = p$ if $p \in P$ is the Purpose Layer actor that authorized the spawn event resulting in $n$'s activation, and $\alpha(n) = \bot$ if and only if $n$ is the root human anchor established at system initialization. The pointer is established exactly once at node provisioning via a Fact Layer atomic write operation and satisfies the following immutability property:
\begin{equation}
\forall n,\ p:\ \bigl(\alpha(n) = p \ \land \ \textit{provisioned}(n)\bigr) \implies \neg\exists\ \textit{modify}(\alpha(n))
\label{eq:parent-pointer-immutable}
\end{equation}
No actor --- including the Purpose Layer after provisioning, the Bounds Engine, any administrative process, or any external principal --- may execute an operation that overwrites, redirects, or deletes $\alpha(n)$ after the provisioning write is confirmed. The Fact Layer write protocol rejects any such attempt at the protocol level, before execution reaches the ledger.
\end{dfn}

The parent pointer is not a reference variable held by the parent node toward its child. It is a Fact Layer property intrinsic to the child node, established at provisioning time and stored in the node's sealed memory envelope. The directionality is architecturally critical: the child carries $\alpha(n)$ as an immutable attribute of itself, not as a mutable external reference held by the parent. This means the chain is traversable from any node outward to its complete lineage regardless of the operational state of any other node in the chain. A parent that has terminated, crashed, or been decommissioned does not affect the child's ability to reference its parent through $\alpha(n)$.

The parent pointer is distinct from a conventional access-controlled pointer in a second critical respect. Access control enforces immutability by restricting which principals may issue modification operations; if sufficient privilege is acquired through credential theft, software vulnerability, or administrative compromise, the access control barrier is circumvented. The Fact Layer write protocol enforces immutability by making the modification operation structurally impossible --- there is no defined mechanism for processing a \textit{modify}($\alpha(n)$) request regardless of any privilege held. There is no privileged override, no emergency bypass, and no administrative role with the capability to alter $\alpha(n)$ after provisioning. This is not a stronger access control policy; it is a different architectural category.

The initialization of $\alpha(n)$ occurs atomically during the spawn event. The Fact Layer records the spawn event in the forensic ledger and writes $\alpha(n) = p$ into the node's sealed memory envelope as part of the same atomic transaction. This separation of authorization (Purpose Layer) from assignment (Fact Layer) is essential: it prevents a parent from claiming to have spawned a node whose pointer was set to a different actor, and it prevents a child from falsifying its own ancestry after provisioning.

\begin{dfn}[Spawn Event]
\label{dfn:spawn-event}
A \emph{spawn event} is the 5-tuple
\begin{equation}
\mathrm{spawn}(c,\ p,\ \sigma_c,\ t,\ \phi)
\label{eq:spawn-event}
\end{equation}
where $c$ is the newly created child node, $p \in P$ is the authorizing Purpose Layer actor, $\sigma_c$ is the authorized execution scope, $t$ is the wall-clock timestamp, and $\phi$ is the cryptographic signature produced by the Purpose Layer over the spawn request. The Fact Layer atomically: (1) creates node $c$, (2) sets $\alpha(c) \leftarrow p$, (3) sets the node's authorized scope to $\sigma_c$, (4) seals the memory envelope, and (5) appends the event to the forensic ledger $\mathcal{L}$.
\end{dfn}

\begin{prop}[Immutability of ParentPointer]
\label{prop:parent-pointer-immutable}
For any child node $c$ with parent $p$, once the spawn event $\mathrm{spawn}(c, p, \sigma_c, t, \phi)$ has been recorded in the forensic ledger $\mathcal{L}$, the value $\alpha(c)$ is fixed for the operational lifetime of $c$. No subsequent operation --- including operations issued by $p$, by any Purpose Layer actor, by the Bounds Engine, or by $c$ itself --- can alter $\alpha(c)$.
\end{prop}

This property distinguishes ParentPointer from conventional mutable parent references in multi-agent systems. A mutable reference can be overwritten by a subsequent administrative operation; an immutable pointer cannot. The distinction is not behavioral --- it is architectural.

% -------------------------------------------------------------------------- %
\subsection{The Chain Operator and Root Human Anchor}
\label{sec:chain-definition}

The chain operator $\Chain{\cdot}$ constructs the complete delegation lineage from any node to the human anchor by recursively applying $\alpha$. It is the primary tool for chain verification, audit, and accountability attribution.

\begin{dfn}[Chain Operator]
\label{def:chain-operator}
For any agent node $a \in \mathcal{N}$ in a Recursive Spawning chain, the \emph{chain} $\Chain{a}$ is defined recursively as:
\begin{equation}
\Chain{a} =
\begin{cases}
\{\, a\,\} \cup \Chain{\alpha(a)} & \text{if } \alpha(a) \neq \bot \\[10pt]
\{\, a,\ h(a) \,\} & \text{if } \alpha(a) = \bot \quad \text{(root human anchor)}
\end{cases}
\label{eq:chain-operator}
\end{equation}

Equivalently, expanding the recursion:
\begin{equation}
\Chain{a} = \bigl\{\, a,\ \alpha(a),\ \alpha^{(2)}(a),\ \alpha^{(3)}(a),\ \ldots,\ n_0 \hmannchor{n_0} \,\bigr\}
\label{eq:chain-expanded}
\end{equation}
where $n_0$ is the unique root satisfying $\alpha(n_0) = \bot$ and $\hmannchor{n_0}$, and $\alpha^{(j)}(a)$ denotes the $j$-fold application of $\alpha$.
\end{dfn}

The chain $\Chain{a}$ is the ordered set of all nodes on the delegation path from $a$ to and including the human anchor. It includes the terminal human anchor node, which is the only node in any RSP chain with $\alpha(n) = \bot$. The chain is always finite: by the RSP termination property, every chain has depth at most $D_{\max}$, so the recursion terminates within at most $D_{\max} + 1$ applications of $\alpha$.

\medskip
\noindent\textbf{Base case: root human anchor.} The root of every RSP chain is a Purpose Layer entity with a human principal designation $\hmannchor{\cdot}$, established at system initialization when a human principal $h$ authorizes the first agent node $n_0$. The Purpose Layer records $h(n_0) = h$ and $\alpha(n_0) = \bot$ in the Fact Layer authorization ledger. This is the only valid base case: the root human anchor has no parent by definition. No machine actor may serve as a chain root in RSP; this is structurally enforced by the Fact Layer pre-commit hook, which rejects any spawn event where the parent node does not satisfy the human anchor reachability predicate.

\medskip
\noindent\textbf{Properties of the chain operator.}
\begin{itemize}
    \item[(C1) \textbf{Determinism}] For any node $a$, $\Chain{a}$ contains exactly one element for each ancestor on the delegation path. There is no alternative chain, no ambiguous parentage, and no non-deterministic branching. $\alpha$ is a total function with unique values, so the chain is uniquely determined.
    \item[(C2) \textbf{Stability}] The immutability of $\alpha$ guarantees that $\Chain{a}$ at any runtime point $t$ is identical to $\Chain{a}$ at $a$'s provisioning moment. The chain is not reconstructed from mutable logs --- it is directly computable by iterating $\alpha$ from any node until reaching $\bot$.
    \item[(C3) \textbf{Compositionality}] For any two nodes $a$ and $b$ in the same RSP chain, either $\Chain{a} \subseteq \Chain{b}$ or $\Chain{b} \subseteq \Chain{a}$. The chains are nested rather than intersecting. No cross-links or merging paths exist in a valid RSP chain, which follows directly from the uniqueness constraint on parent pointers.
\end{itemize}

% -------------------------------------------------------------------------- %
\subsection{Immutability Enforcement: Fact Layer Write Protocol}
\label{sec:fact-layer-enforcement}

The immutability of the parent pointer is enforced by the Fact Layer write protocol --- the \bxthree{} mechanism for making certain state transitions structurally irreversible. The write protocol defines which operations are valid for each node property at each node state; it is not access control with additional layers, but a protocol-level constraint on the state space itself.

\medskip
\noindent\textbf{Write protocol definition.} For the parent pointer $\alpha(n)$, the write protocol defines exactly one valid write operation: the initial provisioning write that sets $\alpha(n) = p$ at node creation. The protocol defines \emph{no valid write operations} for modifying $\alpha(n)$ after provisioning. Any operation --- from any source, with any justification --- that attempts to overwrite, redirect, or delete $\alpha(n)$ after provisioning is rejected at the protocol level. The Fact Layer does not process the request, does not invoke an override pathway, and does not log the attempt as a rejected operation with the option to retry. The request is malformed with respect to the write protocol and is discarded before reaching any evaluation stage.

The rejection is unconditional and source-equivalent:

\begin{itemize}
    \item[\textbf{Failure resilience.}] If the system is in a degraded operational state --- partial network partition, Bounds Engine timeout, Purpose Layer unavailability --- the Fact Layer write protocol for $\alpha(n)$ immutability remains in force. The modification request is rejected regardless of system state, because the write protocol is a local constraint at the Fact Layer, not a distributed agreement protocol that could be blocked by network failure.
    \item[\textbf{Administrative override absence.}] There is no administrative role, system operator, or security principal with the capability to modify $\alpha(n)$. This is a deliberate architectural omission. In conventional systems, administrative compromise leads directly to immutability violation because the administrative role holds override authority. In RSP, there is no override authority for $\alpha(n)$ by design.
    \item[\textbf{Source-equivalence.}] The Fact Layer write protocol treats all modification requests identically regardless of source. A request from the Purpose Layer, the Bounds Engine, an authenticated administrator, or an external adversarial actor receives the same response: protocol-level rejection. There is no privileged source that bypasses this check.
    \item[\textbf{Atomic warrant-pointer creation.}] The parent pointer and its Purpose Layer warrant are created as a single atomic Fact Layer write. There is no temporal window in which $\alpha(n)$ is established without a matching warrant, or in which a warrant exists without a corresponding parent pointer.
\end{itemize}

\medskip
\noindent\textbf{The sealed memory envelope.} The immutability guarantee is reinforced by the sealed memory envelope mechanism:

\begin{enumerate}
    \item \textbf{Encrypted at rest.} $\alpha(n)$ is encrypted in the node's persistent memory envelope using a Fact Layer-only symmetric key. Even a compromised Bounds Engine cannot read $\alpha(n)$ because it never has access to the decryption key.
    \item \textbf{HMAC-signed.} The envelope carries an HMAC-SHA-256 signature computed over its contents using a Fact Layer-only key. Any modification to the ciphertext invalidates the signature.
    \item \textbf{Decrypted only for reads.} The Fact Layer decrypts the envelope only when reading $\alpha(n)$ to service a chain-traversal or audit request. The plaintext is never exposed to the Bounds Engine or the child node's own code.
    \item \textbf{Re-sealed after reads.} After each authorized read, the Fact Layer re-encrypts and re-signs the envelope, preventing replay attacks based on stale plaintext.
\end{enumerate}

The Fact Layer operates as a standalone, isolated process with its own memory space, inaccessible to the child node's address space. The combination of the protocol-level write constraint and the sealed envelope ensures that $\alpha(n)$ is modification-resistant in all deployment contexts.

\medskip
\noindent\textbf{Why access control is insufficient.} The distinction between access-control-based immutability and protocol-level immutability is the distinction between a promise and a guarantee. In an access-control model, immutability is enforced by restricting which principals may issue modification operations; if an attacker acquires sufficient privileges through credential theft, insider compromise, or software vulnerability, the barrier is circumvented and immutability is violated. The immutability was a policy, enforceable until it was not.

In the Fact Layer write protocol, the immutability of $\alpha(n)$ is a property of the protocol itself. There is no operation defined for modifying $\alpha(n)$ after provisioning, so no amount of privilege elevation, credential compromise, or software vulnerability can produce a valid modification operation. The protocol does not enforce immutability --- it makes immutability the only valid state transition. The distinction is architectural, not a matter of degree.

This is the precise failure mode the Fact Layer write protocol eliminates: retroactive manipulation of chain topology for the purpose of accountability laundering. Without protocol-level immutability, an adversarial actor could re-parent a node to a favorable parent after the fact, obscuring the original authorization chain and making drift detection impossible. With protocol-level immutability, this manipulation is structurally impossible.

% -------------------------------------------------------------------------- %
\subsection{Spawn Authorization: Purpose Layer Role}
\label{sec:purpose-layer-authorization}

The Purpose Layer plays two structurally necessary roles in the parent pointer lifecycle: it originates the spawn authorization policy that defines conditions for creating a new parent pointer, and it issues the spawn warrant that is the prerequisite for the Fact Layer provisioning write. These roles are inseparable from the parent pointer's immutability: if the Purpose Layer could issue warrants after the fact, or if a machine actor could create a parent pointer without a Purpose Layer warrant, the chain's integrity guarantees would be void.

\medskip
\noindent\textbf{Authorization policy origination.} At system initialization, the Purpose Layer defines the spawn authorization policy $P_{\text{spawn}}$ and records it in the Fact Layer authorization ledger. $P_{\text{spawn}}$ specifies the mandatory conditions for any valid spawn event:

\begin{enumerate}
    \item[(S1) \textbf{Scope refinement.}] The child scope must be a strict subset of the parent scope: $R(n_{\text{child}}) \subset R(n_{\text{parent}})$. No lateral expansion, no superset, no equality. Each spawn event narrows the delegation scope, preserving the intent-refinement property of the chain.
    \item[(S2) \textbf{Operational necessity.}] The parent must declare, in the Fact Layer authorization ledger, the specific operational condition that necessitates spawning a child and cannot be resolved within the parent's scope.
    \item[(S3) \textbf{Minimum scope proportionality.}] The child scope must be the minimum scope necessary to address the declared operational condition. Broader scopes are not authorized.
    \item[(S4) \textbf{Depth bound compliance.}] The proposed spawn must satisfy $\text{depth}(n_{\text{parent}}) + 1 \leq D_{\max}$. If this condition fails, the spawn is rejected and mandatory escalation is triggered.
\end{enumerate}

These conditions are non-negotiable. The Purpose Layer does not issue a spawn warrant if any condition is unmet. The Fact Layer pre-commit hook independently re-verifies these conditions before recording the provisioning write.

\medskip
\noindent\textbf{Spawn warrant issuance.} When all four conditions are satisfied, the Purpose Layer issues a spawn warrant $w$ --- a cryptographically signed, Fact Layer-recorded authorization record:

\begin{equation}
w = \bigl\langle \text{parent}=n_i,\ \text{child}=n_{i+1},\ R(n_{i+1}),\ \text{justification},\ \text{timestamp},\ \sigma_{\text{Purpose Layer}} \bigr\rangle
\label{eq:spawn-warrant-structure}
\end{equation}

The warrant is written to the Fact Layer authorization ledger atomically with the provisioning write that establishes $\alpha(n_{i+1}) = n_i$. The atomicity is critical: the parent pointer and the spawn warrant are created as a single indivisible ledger entry.

\medskip
\noindent\textbf{Sole authorization origin.} No machine actor may establish a parent pointer without a Purpose Layer warrant. This is a structural consequence of the Fact Layer write protocol: the pre-commit hook verifies the existence of a valid Purpose Layer warrant before recording any provisioning write. A request to establish $\alpha(n)$ without a matching warrant is rejected at the pre-commit hook, before the operation reaches the ledger. The Purpose Layer's exclusive role as authorization origin prevents the chain from being bootstrapped by a machine actor acting autonomously.

% -------------------------------------------------------------------------- %
\subsection{Chain Traversal Algorithm}
\label{sec:chain-traversal}

The chain traversal algorithm is the runtime procedure for constructing the delegation chain from any node to its human anchor. It is used by the Bailout Protocol for exception escalation, by auditors for accountability verification, and by the Fact Layer's continuous orphan detection.

\medskip
\noindent\textbf{Algorithm: Chain-Traverse}

\begin{algorithm}
\caption{Chain-Traverse($n$) --- Build the delegation chain from node $n$ to the human anchor}
\label{alg:chain-traverse}
\begin{algorithmic}[1]
\REQUIRE $n$ is a provisioned RSP node with $\alpha$ defined
\ENSURE $\Chain{n}$ --- the ordered set of all ancestor nodes including the human anchor
\STATE $chain \leftarrow \langle\ \rangle$ \COMMENT{Initialize empty ordered list}
\STATE $current \leftarrow n$ \COMMENT{Start at the query node}
\WHILE{$\text{true}$}
    \STATE $\text{append}(chain,\ current)$ \COMMENT{Add current node to chain}
    \IF{$\alpha(current) = \bot$}
        \STATE \textbf{break} \COMMENT{Root human anchor reached --- chain complete}
    \ENDIF
    \STATE $current \leftarrow \alpha(current)$ \COMMENT{Move to parent}
\ENDWHILE
\RETURN $chain$
\end{algorithmic}
\end{algorithm}

\medskip
\noindent\textbf{Correctness arguments.}

\textit{Termination.} The algorithm terminates because RSP chains are depth-bounded. Each iteration of the while-loop moves $current$ to $\alpha(current)$, which is strictly closer to the root. Since the chain has finite depth at most $D_{\max} + 1$, the loop executes at most $D_{\max} + 1$ iterations before $\alpha(current) = \bot$ is reached. The loop therefore terminates in finite time regardless of the operational state of any node in the chain.

\textit{Correctness of the returned set.} By Definition~\ref{def:chain-operator}, $\Chain{a}$ is the set containing $a$, $\alpha(a)$, $\alpha(\alpha(a))$, and so on, until the root. The algorithm produces exactly this sequence: it starts at $a$, appends each node, moves to its parent via $\alpha$, and stops when $\alpha(current) = \bot$ is reached. The output is identical to $\Chain{a}$ by construction.

\textit{Uniqueness of result.} The algorithm produces a unique result for any input node because $\alpha$ is a total function (defined for every provisioned node) and injective on the chain direction (every node has exactly one parent, and the root has no parent). There is no branching, no alternative path, and no non-deterministic choice at any step.

\textit{Immutability guarantee.} The algorithm reads $\alpha$ values but never writes them. The traversal is read-only and does not interact with the Fact Layer write protocol. The algorithm can be executed at any runtime point without affecting the immutability of the chain it traverses, and without creating any modification opportunities for adversarial actors.

\medskip
\noindent\textbf{Algorithm: Chain-Verify}

The chain traversal algorithm enables Chain-Verify, which confirms the structural validity of a delegation chain at any runtime point:

\begin{algorithm}
\caption{Chain-Verify($n$) --- Verify the authorization chain from node $n$ to human anchor}
\label{alg:chain-verify}
\begin{algorithmic}[1]
\REQUIRE $n$ is a provisioned RSP node
\ENSURE $\text{Boolean}$ indicating whether the chain is a structurally valid RSP chain
\STATE $chain \leftarrow \text{Chain-Traverse}(n)$
\STATE $m \leftarrow \text{len}(chain)$

\STATE \COMMENT{\textit{Step V1: Verify chain terminates at human anchor}}
\IF{$\alpha(chain[m-1]) \neq \bot$}
    \STATE \RETURN $\text{false}$ \COMMENT{Root is not $\bot$ --- malformed chain}
\ENDIF
\IF{$\neg h(chain[m-1])$}
    \STATE \RETURN $\text{false}$ \COMMENT{Root is not a designated human principal}
\ENDIF

\STATE \COMMENT{\textit{Step V2: Verify each spawn event has a Purpose Layer warrant}}
\FOR{$i \leftarrow 0$ \TO $m-2$}
    \STATE $child \leftarrow chain[i]$
    \STATE $parent \leftarrow chain[i+1]$
    \IF{$\neg \text{warrant\_exists}(child,\ parent)$}
        \STATE \RETURN $\text{false}$ \COMMENT{Missing Purpose Layer warrant for spawn}
    \ENDIF
\ENDFOR

\STATE \COMMENT{\textit{Step V3: Verify scope refinement at each link}}
\FOR{$i \leftarrow 0$ \TO $m-2$}
    \STATE $child \leftarrow chain[i]$
    \STATE $parent \leftarrow chain[i+1]$
    \IF{$\neg (R(child) \subset R(parent))$}
        \STATE \RETURN $\text{false}$ \COMMENT{Scope refinement constraint violated}
    \ENDIF
\ENDFOR

\STATE \COMMENT{\textit{Step V4: Verify depth bound compliance at each node}}
\FOR{$i \leftarrow 0$ \TO $m-1$}
    \IF{$\text{depth}(chain[i]) > D_{\max}$}
        \STATE \RETURN $\text{false}$ \COMMENT{Depth bound exceeded at node}
    \ENDIF
\ENDFOR

\RETURN $\text{true}$
\end{algorithmic}
\end{algorithm}

Chain-Verify returns \texttt{true} if and only if all four verification conditions hold simultaneously. V1 confirms the chain terminates at a human anchor; V2 confirms every spawn was authorized by the Purpose Layer; V3 confirms scope refinement was respected at every link; V4 confirms the depth bound was respected throughout. A chain that passes Chain-Verify is, by the RSP definitions, a fully authorized delegation chain terminating at a human principal.

The algorithm is deterministic and produces the same result regardless of which node initiates verification. There is no self-serving interpretation of authorization that a drifted or orphaned node could produce, because Chain-Verify inspects the complete chain topology and warrants in the Fact Layer ledger, not the node's own claims about its authorization.

\medskip
\noindent\textbf{Constant-time escalation bound.} The depth bound $D_{\max}$ provides a constant-time upper bound on chain traversal cost. The worst-case number of pointer dereferences required to reach the human anchor from any node is $D_{\max} + 1$, independent of the total number of nodes in the system. A system with ten thousand spawned nodes and $D_{\max} = 5$ requires at most six pointer dereferences to reach a human --- the same as a system with ten nodes. Mutable parent pointers would allow the chain to be extended retroactively, making the effective depth potentially unbounded and the worst-case escalation time unbounded as well. Immutability is what makes the depth bound a structural guarantee, not merely a policy.

In the Bailout Protocol context, chain traversal is used to identify the correct Human Accountability Anchor when an agent at any depth encounters an unresolvable exception. The upward propagation of the exception signal follows the same parent-pointer structure as chain traversal, but uses the asynchronous trigger pathway rather than the synchronous verification pathway. Chain traversal ensures that this walk never diverges, never shortcuts, and never fails to reach a human.

\medskip
\noindent\textbf{Computational cost.} Chain-Traverse runs in $O(d)$ time where $d$ is the chain depth (at most $D_{\max} + 1$). Chain-Verify runs in $O(m)$ for the traversal plus $O(m)$ for each of the three verification loops, yielding $O(m)$ total where $m$ is chain length, also bounded by $D_{\max} + 1$. Both algorithms are linear in chain length and do not depend on the total number of nodes in the system.

% -------------------------------------------------------------------------- %
\subsection{Summary: Immutability as Architectural Constraint}
\label{sec:immutability-summary}

The immutable parent pointer is the load-bearing constraint of the Recursive Spawning Protocol. Without it, the chain is mutable --- subject to retroactive modification by actors with sufficient privilege --- and the RSP's correctness properties (drift prevention, chain integrity, termination) collapse to policy conventions rather than structural guarantees.

The immutability is architectural --- not a stronger access control policy --- because it is enforced by the Fact Layer write protocol, which defines no valid modification operation for $\alpha(n)$ after provisioning. This is distinct from access-control-based immutability in four ways: there is no override path for any actor; the immutability holds under all system failure conditions including those that disable other enforcement mechanisms; all sources receive equivalent enforcement treatment regardless of privilege level; and the parent pointer and its Purpose Layer warrant are created atomically with no temporal window for inconsistency.

The chain operator $\Chain{a}$ is uniquely determined by the immutable parent pointers in the chain. The chain traversal algorithm computes this operator correctly and terminates. Chain-Verify provides the runtime verification procedure for chain structural validity. Together, these mechanisms make the delegation chain an immutable, verifiable, and auditable runtime artifact --- the foundation on which all other RSP correctness properties rest.
